\documentclass[10pt,a4paper]{amsart}
\usepackage[margin=19mm]{geometry}
\usepackage{amsmath,amssymb,mathtools}
\usepackage[scr=rsfs,cal=euler]{mathalfa}
\usepackage{xfrac}
\usepackage[unicode,hidelinks,bookmarks=false]{hyperref}
\newcommand{\ie}{i.e.\ }
\newcommand{\cf}{cf.\ }
\newcommand{\resp}{respectively\ }
\newcommand{\Subsection}{\subsection}
\providecommand{\nobreakdash}{\nobreak\hbox{-}}
\setlength{\emergencystretch}{2em}
\multlinegap0pt
\usepackage{amssymb}
\usepackage{dsfont}
\usepackage[shortlabels]{enumitem}
\usepackage{cancel}
\usepackage{xcolor}
\usepackage{tikz}
\newtheorem{theorem}{Theorem}
\newtheorem{lemma}{Lemma}
\newcommand{\N}{\mathbb N}
\newcommand{\R}{\mathbb R}
\newcommand{\bq}{\begin{equation}}
\newcommand{\calA}{\mathcal A}
\newcommand{\calD}{\mathcal D}
\newcommand{\calK}{\mathcal K}
\newcommand{\calM}{\mathcal M}
\newcommand{\calP}{\mathcal P}
\newcommand{\calR}{\mathcal R}
\newcommand{\diam}{\text{diam}}
\newcommand{\dr}{\textnormal{d}r}
\newcommand{\ds}{\textnormal{d}s}
\newcommand{\dt}{\textnormal{d}t}
\newcommand{\dx}{\textnormal{d}x}
\newcommand{\dy}{\textnormal{d}y}
\newcommand{\dz}{\textnormal{d}z}
\newcommand{\e}{\varepsilon}
\newcommand{\eq}{\end{equation}}
\DeclareMathOperator*{\esssup}{ess\,sup}
\renewcommand{\ge}{\geqslant}
\renewcommand{\geq}{\geqslant}
\newcommand{\intr}{\int_{\R^d}}
\newcommand{\intrr}{\iint_{\R^d \times \R^d}}
\renewcommand{\le}{\leqslant}
\newcommand{\lt}{\left}
\newcommand{\pa}{\partial}
\newcommand{\rd}{\textnormal{d}}
\newcommand{\rt}{\right}
\newcommand{\supp}{\text{supp}}
\title{Lagrangian formulation and Eulerian closure\\in alignment dynamics}
\author{Jos\'e A. Carrillo, Young-Pil Choi, Eitan Tadmor}
\date{Source: arXiv:2604.10253v1 (CC BY 4.0)}
\begin{document}
\maketitle

\noindent\emph{Source and licence.} Lagrangian formulation and Eulerian closure\\in alignment dynamics. Version arXiv:2604.10253v1 (CC BY 4.0), distributed by arXiv under the Creative Commons Attribution 4.0 International licence, http://creativecommons.org/licenses/by/4.0/, as recorded in the arXiv licence metadata for this exact version and shown on the abstract page https://arxiv.org/abs/2604.10253v1. Full licence notice: \texttt{licenses/arxiv-2604.10253-CC-BY-4.0.txt}.\par\smallskip

\noindent\textit{Editorial scope.} Counted unit: the article's theorem thm:asymp\_closure (source0.tex lines 671--714). The article's complete original proof of it is delivered with it (source0.tex lines 1997--2233, one \texttt{proof} environment of 237 lines, which the article places away from the statement). No mathematics is written, repaired or re-derived: the statement and the proof are the article's own lines, in the article's own notation, and no retained line is substituted or re-flowed.  Retained uncounted as the source-local context the proof needs: lines 311--331; lines 320--325; lines 369--374; lines 423--438; lines 551--594; lines 1266--1293; lines 1305--1313; lines 1466--1479; lines 1911--1925; lines 3003--3003.  Nothing was omitted from any retained span, so the package carries no omission record.  No retained line contains a citation, so the wrapper carries no bibliography and no bibliography entry.  No retained line contains a figure environment, an image-inclusion command or drawing code, so no image is delivered and the package declares no asset dependency.  Editorial formatting only, in addition: an AMS \texttt{amsart} preamble that restates the article's own declarations and macros used by the retained lines, namely the environments \texttt{theorem}, \texttt{lemma} on separate counters, together with the standard packages those lines need.  The retained environments print in this document as Theorem 1, Lemma 1 in order of appearance, whereas the article prints the same items with the numbering of its own full document.  The complete native source of the article as distributed in the arXiv e-print for this exact version is bundled unchanged as \texttt{sources/198254/source0.tex}.
%%%%%%%%%%%%%%%%%%%%%%%%%%%%%%%%%%%%%%%%%%%%%%%%%%%%%%%%%%%%%%%%%%%%%%%%%%%%%%%%%5
%
%
%                        Section: Introduction 
%
%
%%%%%%%%%%%%%%%%%%%%%%%%%%%%%%%%%%%%%%%%%%%%%%%%%%%%%%%%%%%%%%%%%%%%%%%%%%%%%%%%%
\section{Introduction}%\label{sec:intro}
We consider the following infinite-dimensional, measure-dependent system of ordinary differential equations describing the Lagrangian evolution of a continuum ensemble of interacting agents:
\begin{align}\label{Lag_p}
\begin{aligned}
\pa_t \eta_t &= v_t, \quad (t,x) \in \R_+ \times \supp(\rho_0),\cr
\pa_t v_t &=\kappa \intr \phi(\eta_t(x) - \eta_t(y)) G_p(v_t(y) - v_t(x))\,\rho_0(\dy), \quad  G_p(\xi):=|\xi|^{p-2}\xi,
\end{aligned}
\end{align}
which we call the \emph{Lagrangian $p$-alignment formulation}. When $p=2$, the interaction becomes linear in the velocity difference, and we refer to \eqref{Lag_p} as the \emph{Lagrangian alignment formulation}. Here $\kappa>0$ denotes the coupling strength, $\rho_0 \in \calP(\R^d)$ is a given probability measure describing the initial spatial distribution of particles, and the initial data are prescribed as
\bq\label{ini:Lag_p}
(\eta_t(x), v_t(x))|_{t=0} = (x, u_0(x)), \quad x \in \supp(\rho_0).
\eq
Throughout the paper, we assume that the communication kernel $\phi \in C^1 \cap W^{1,\infty}(\R^d;\R_+)$ is radially symmetric and non-increasing. In particular, $\phi(z)=\phi(|z|)$ and $\phi$ is even, ensuring symmetry of
pairwise interactions.

\begin{align}
\begin{aligned}
\pa_t \eta_t &= v_t, \quad (t,x) \in \R_+ \times \supp(\rho_0),\cr
\pa_t v_t &=\kappa \intr \phi(\eta_t(x) - \eta_t(y)) G_p(v_t(y) - v_t(x))\,\rho_0(\dy), \quad  G_p(\xi):=|\xi|^{p-2}\xi,
\end{aligned}
\end{align}

\begin{align}\label{EA}
\begin{aligned}
&\pa_t \rho_t + \nabla \cdot (\rho_t u_t) = 0, \quad t>0, \ x \in \R^d, \cr
&\pa_t (\rho_t u_t) + \nabla \cdot (\rho_t u_t \otimes u_t) = \rho_t\calA_p[\rho_t,u_t],
\end{aligned}
\end{align}

 \begin{theorem}\label{thm:ext} Let  $p\ge2$ and suppose that $\phi:\R^d\to[0,\infty)$ is bounded and Lipschitz continuous. 
For any $u_0\in  L^\infty(\rho_0; \R^d)$, the Lagrangian $p$-alignment system \eqref{Lag_p}--\eqref{ini:Lag_p} admits a unique global solution
\bq\label{reg_Lag}
(\eta - {\rm id},v)\in C^1([0,\infty);L^\infty(\rho_0)\times  L^\infty(\rho_0))
\eq
satisfying
\[
\|v_t\|_{L^\infty(\rho_0)} \le \|u_0\|_{L^\infty(\rho_0)},  \quad  \|\eta_t-{\rm id}\|_{L^\infty(\rho_0)} \le t \|u_0\|_{L^\infty(\rho_0)} \quad\text{for all }t\ge0.
\]
 Moreover, assume that $\rho_0$ is compactly supported and the communication kernel $\phi$ is ``heavy-tailed'' in the sense that
\bq\label{condi_phi}
\int \limits^\infty \phi(r)\,\dr = \infty.
\eq
Then the velocity diameter asymptotically vanish, i.e.,
$\rd_v(t) \to 0 \  \text{as } t \to \infty$.
 \end{theorem}

\begin{theorem} \label{thm:ERA}
Assume $p \geq 2$ and the hypotheses of Theorem \ref{thm:ext}. Let $(\rho_t, u_t, \tau_t)$ be the Eulerian variables induced by the Lagrangian solution $(\eta_t,v_t)$ through pushforward and disintegration as described above. Then $(\rho_t, u_t, \tau_t)$ is a global solution to 
\begin{align}\label{eq:ERA}
\begin{aligned}
\pa_t\rho_t+\nabla \cdot(\rho_t u_t)&=0, \\
\pa_t(\rho_t u_t)+\nabla \cdot(\rho_t u_t\otimes u_t + \tau_t)
&=\rho_t\lt(\calA_p[\rho_t,u_t]+\calR_p[\rho_t,u_t]\rt), 
\end{aligned}
\end{align}
in the sense of distributions on $(0,\infty)\times\R^d$, with initial data 
\[
(\rho_t, u_t, \tau_t)|_{t=0} = (\rho_0, u_0, 0),
\]
and satisfying\footnote{The continuity statements 
\[
\rho\in C([0,\infty);\calP(\R^d)),\quad
m\in C([0,\infty);\calM(\R^d;\R^d)) 
\]
refer to the narrow topology on $\calP(\R^d)$ and the weak$^\ast$
topology on the indicated spaces of Radon measures. In particular, a strong time-regularity of $u_t$ or $\theta_t$ is not implied.
}
\[
\rho \in C([0,\infty);\calP(\R^d)),\quad m=\rho u \in C([0,\infty);\calM(\R^d;\R^d)), \quad \tau \in L^\infty_{\rm loc}([0,\infty);\calM(\R^d;\R^{d \times d})).
\]
Here, the nonlinear defect force is
\[
\calR_p[\rho_t,u_t](z) :=\kappa\intr\phi(z-\zeta) \calK_t(z,\zeta)\,\rho_t(\rd\zeta),
\]
where, denoting by $\{\nu_{t,z}\}$ any disintegration of $\rho_0$ along $\eta_t$, and $\omega_t(x):=v_t(x)-u_t(\eta_t(x))$ which has zero fibre mean, i.e., $\intr \omega_t\,\nu_{t,z}(\dx)=0$, we set
\bq\label{calK}
\calK_t(z,\zeta) := \intrr   G_p(u_t(\zeta)-u_t(z)+\omega_t(y)-\omega_t(x)) \,\nu_{t,\zeta}(\dy) \nu_{t,z}(\dx) - G_p(u_t(\zeta)-u_t(z)).
\eq

Moreover, $(\rho,u,\tau)$ satisfies the global energy inequality:
for all $t\geq 0$,
\begin{align}\label{ERA_energy}
\begin{aligned}
&\frac12 \intr \left(|u_t|^2+{\rm tr}\,\theta_t(z)\right)\rho_t(\dz)\cr
&\quad +\frac{\kappa}{2} \int_0^t\iiiint_{\R^d \times \R^d \times \R^d \times \R^d} \phi(z-\zeta)\,|v_s(y)-v_s(x)|^p \,\nu_{s,z}(\dx)\nu_{s,\zeta}(\dy) \rho_s(\dz)\rho_s(\rd\zeta)\, \ds \cr
&\qquad \le \frac12\intr |u_0|^2 \,\rho_0(\dz).
\end{aligned}
\end{align}
In the class of Eulerian triples arising from the Lagrangian flow of Theorem \ref{thm:ext}, this solution is unique.
\end{theorem}

We now establish an $L^\infty$ maximum principle for $v$ using the upper Dini derivative. Fix a unit vector $e\in\R^d$ and set $s_t(x):=v_t(x)\cdot e$. Then, $s$ satisfies
\[
\pa_t s_t(x) =\intr K_t(x,y) (s_t(y)-s_t(x))\,\rho_0(\dy),
\]
where
\[
K_t(x,y):=\kappa \phi ((x+w_t(x))-(y+w_t(y)))|v_t(y)-v_t(x)|^{p-2}.
\]
Since $\kappa\ge0$ and $\phi\ge0$, we have $K_t\ge0$. For each fixed $t$, local existence guarantees  $\|v_t\|_{L^\infty(\rho_0)}<\infty$, and thus
\[
\|K_t\|_{L^\infty(\rho_0\otimes\rho_0)} \le \kappa \|\phi\|_{L^\infty} \lt(2\|v_t\|_{L^\infty(\rho_0)}\rt)^{p-2}<\infty.
\]
Define $M_e(t):=\esssup_{x}s_t(x)$, and let 
\[
D^+M_e(t):=\limsup_{h\downarrow0}\frac{M_e(t+h)-M_e(t)}{h}
\]
be its upper right Dini derivative. Fix $t\ge0$ and $\varepsilon>0$, and then choose $x_\e$ such that $s_t(x_\e)\ge M_e(t)-\e$. Then 
\[
s_t(y)-s_t(x_\e) \le \e \quad \text{for $\rho_0$-a.e.  $y$}, 
\]
and hence,
\[
\pa_t s_t(x_\e) =\intr K_t(x_\e,y)(s_t(y)-s_t(x_\e))\,\rho_0(\dy) \le \e \|K_t\|_{L^\infty(\rho_0\otimes\rho_0)}.
\]
By the standard Dini envelope argument (letting $\varepsilon\downarrow0$ at fixed $t$), we conclude $D^+M_e(t)\le0$ for all $t\geq 0$ and unit vector $e$.   Hence, we have
\bq\label{eq:v-Linf-decay}
\|v_t\|_{L^\infty(\rho_0)} =\sup_{|e|=1} M_e(t)\le \sup_{|e|=1} M_e(0)\le \|u_0\|_{L^\infty(\rho_0)} \quad\text{for all }t<T_{\max}.
\eq

\begin{lemma}\label{lem:apr} Let $(\eta,v)$ be a global classical solution to the system \eqref{Lag_p}. Then we have
\[
\frac{\rd}{\dt} \intr \eta_t(x)\,\rho_0(\dx) = \intr v_t(x)\,\rho_0(\dx), \quad \frac{\rd}{\dt}\intr v_t(x)\,\rho_0(\dx) =0,
\]
and
\[
\frac12 \frac{\rd}{\dt} \intr |v_t(x)|^2\,\rho_0(\dx) + \frac\kappa2\intrr \phi(\eta_t(x)-\eta_t(y))|v_t(x) - v_t(y)|^p \,\rho_0(\dx) \rho_0(\dy) =0.
\]
\end{lemma}

\begin{theorem}\label{thm:disint}
Let $X$, $\bar X$ be Radon separable metric spaces, $\mu \in \calP(X)$, let $\pi: X \to \bar X$ be a Borel map and let $\nu = \pi_\# \mu \in \calP(\bar X)$. Then there exists a $\nu$-a.e. uniquely determined Borel family of probability measures $\{\mu_{\bar x}\}_{\bar x \in \bar X} \subset \calP(X)$ such that
\[
\mu_{\bar x} (X \setminus \pi^{-1}(\{\bar x\})) = 0 \quad \text{for $\nu$-a.e. $\bar x \in \bar X$},
\]
and for every Borel map $f: X \to [0,+\infty)$,
\[
\int_X f(x)\,\mu(\dx) = \int_{\bar X} \lt(  \int_{\pi^{-1}(\bar x)} f(x) \,\mu_{\bar x}(\dx)\rt)  \nu(\rd\bar x).
\]
In particular, if $X = X_1 \times X_2$, $\bar X = X_1$, and $\mu \in \calP(X_1 \times X_2)$ with first marginal $\nu = \pi_1 {}_\# \mu$, then one can identify each fibre $\pi_1^{-1}(x_1) = \{x_1\} \times X_2$ with $X_2$ and find a Borel family $\{\mu_{x_1}\}_{x_1 \in X_1} \subset \calP(X_2)$ (unique $\nu$-a.e.) such that
\[
\mu = \int_{X_1} \mu_{x_1} \nu(\rd x_1).
\]
\end{theorem} 

\begin{lemma}\label{lem:asymp_closure}
Let $p \geq 2$, assume the hypotheses of Theorem \ref{thm:ext}, and let $(\rho_t,u_t,\tau_t)$ be the Eulerian objects induced by the Lagrangian solution as in Theorem \ref{thm:ERA}. Then the Reynolds stress vanishes asymptotically in the sense that
\[
\|\tau_t\|_{\calM(\R^d;\R^{d\times d})} \to 0 \quad \text{as }t\to\infty.
\]
In addition, when $p>2$, the nonlinear defect force also vanishes asymptotically:
\[
\|\rho_t \calR_p[\rho_t,u_t]\|_{\calM(\R^d;\R^d)} \to 0 \quad \text{as }t\to\infty.
\]
Moreover, the Eulerian velocity diameter decays to zero in the essential sense:
\[
\esssup_{(z,\zeta)\in S_t\times S_t}|u_t(z)-u_t(\zeta)|
\to 0 \quad \text{as }t\to\infty, \qquad S_t=\text{supp}\,\rho_t.
\]
\end{lemma}

\section{Kinetic formulation associated with the Lagrangian flow}\label{app:kinetic}

\begin{theorem} \label{thm:asymp_closure}
Let $(\rho_t,u_t,\tau_t)$ be the Eulerian fields associated with a global  Lagrangian solution constructed in Theorem \ref{thm:ext}.  Assume that $\rho_0$ is compactly supported and the communication kernel satisfies the heavy-tail condition \eqref{condi_phi}. Then the following hold.

\medskip
\noindent
\textbf{(Asymptotic closure)} The Reynolds stress vanishes asymptotically:
\[
\|\tau_t\|_{\calM(\R^d;\R^{d\times d})}\to0 \quad\text{as }t\to\infty.
\]
In addition, when $p>2$, the nonlinear defect force also vanishes asymptotically:
\[
\|\rho_t \calR_p[\rho_t,u_t]\|_{\calM(\R^d;\R^d)}\to0 \quad\text{as }t\to\infty.
\]
Consequently, the defect terms in the ERA system vanish as $t\to\infty$, and the dynamics become asymptotically consistent with the mono-kinetic Euler--alignment system \eqref{EA}.


\medskip
\noindent
\textbf{(Asymptotic flocking)} The diameter of Eulerian velocities decays to zero in the essential sense:
\[
\esssup_{(z,\zeta)\in S_t\times S_t}|u_t(z)-u_t(\zeta)|
\to 0 \quad \text{as }t\to\infty, \qquad S_t :=\text{supp}\,\rho_t.
\]



\medskip
\noindent
\textbf{(Asymptotic mono-kinetic dynamics)}
Let $T>0$ and let $\mu_t=(\eta_t,v_t)_\#\rho_0$ denote the lifted kinetic measure. Then there exist a sequence $R_n\to\infty$ and a family of probability measures $\{\rho_t^\ast\}_{t\in[0,T]}\subset\calP(\R^d)$ such that
\[
\mu_{R_n+t} \rightharpoonup \rho_t^\ast(\rd z)\,\delta_{\bar u}(\rd\xi) \quad
\text{in }\calM_{\mathrm{loc}}(\R^d\times\R^d)
\]
for every $t\in[0,T]$, where
\[
\bar u=\intr u_0(x)\rho_0(\rd x).
\]
Moreover,
\[
\rho_t^\ast=(z\mapsto z+t\bar u)_\#\rho_0^\ast,
\]
where $\rho_0^\ast\in\mathcal P(\R^d)$ denotes the weak limit of the spatial marginals at the initial shifted time.
\end{theorem}

\begin{proof}[Proof of Theorem \ref{thm:asymp_closure}]
The first assertion is exactly Lemma \ref{lem:asymp_closure}. We therefore only prove the second assertion. We proceed in several steps.

\medskip
\noindent\textbf{Step 1. Time translation and kinetic weak formulation on a fixed window.}
Fix $T>0$. For each $R>0$, introduce the translated lifted measure
\[
\mu^{(R)}_s := (\eta_{R+s},v_{R+s})_\#\rho_0, \quad 0\le s\le T.
\]
By construction, $\mu^{(R)}_s\in\calP(\R^d\times\R^d)$.

As shown in Appendix~\ref{app:kinetic}, the lifted measure $\mu_t$ associated with the Lagrangian dynamics satisfies the kinetic equation
\[
\partial_t\mu_t +\nabla_z\cdot(\xi\mu_t) +\nabla_\xi\cdot\bigl(F_p[\mu_t]\mu_t\bigr)=0 \quad\text{in }\calD'((0,\infty)\times\R^d\times\R^d),
\]
where
\[
F_p[f](z,\xi) = \kappa \iint_{\R^d\times\R^d} \phi(z-z')\,G_p(\xi'-\xi)\,f(\rd z',\rd\xi'),
\quad
G_p(w)=|w|^{p-2}w.
\]

Consequently, the translated measure $\mu^{(R)}$ satisfies the weak formulation
\begin{align}\label{eq:kinetic_weak_R}
\begin{aligned}
&\int_0^T\intrr \lt(-\partial_s\Phi(s,z)\cdot\xi -\nabla_z\Phi(s,z):(\xi\otimes\xi)\rt) \mu^{(R)}_s(\rd z,\rd\xi) \,\rd s \\ 
&\quad = \int_0^T\intrr \Phi(s,z)\cdot F_p[\mu^{(R)}_s](z,\xi) \,\mu^{(R)}_s(\rd z,\rd\xi) \,\rd s
\end{aligned}
\end{align}
for every $\Phi\in C_c^\infty((0,T)\times\R^d;\R^d)$.

In particular, taking test functions depending only on $(s,z)$ yields the spatial continuity equation
\[%\begin{equation}\label{eq:continuity_muR}
\int_0^T\intrr
\lt(-\partial_s\varphi(s,z)-\xi\cdot\nabla_z\varphi(s,z)\rt)  \mu^{(R)}_s(\rd z,\rd\xi)\,\rd s=0
\]%\end{equation}
for all $\varphi\in C_c^\infty((0,T)\times\R^d)$.
 
\medskip
\noindent\textbf{Step 2. Local compactness on the support of the test function.}
Let $K\subset\R^d$ be a compact set such that $\supp\Phi\subset (0,T)\times K$. Since we only test the kinetic formulation \eqref{eq:kinetic_weak_R} against $\Phi$, all space--time integrals involve only $(s,z)\in[0,T]\times K$.

By the $L^\infty$ maximum principle established in the proof of Theorem \ref{thm:ext} (see \eqref{eq:v-Linf-decay}), we have the uniform bound
\[%\bq\label{eq:vel_bound}
\|v_t\|_{L^\infty(\rho_0)} \le \|u_0\|_{L^\infty(\rho_0)}
\quad\text{for all } t\ge0,
\]%\eq
and hence $|v_t(x)|\le M$ for $\rho_0$-a.e.  $x$ with $M:=\|u_0\|_{L^\infty(\rho_0)}$. Consequently, for every $R>0$ and $s\in[0,T]$, the lifted measure
\[
\mu^{(R)}_s=(\eta_{R+s},v_{R+s})_\#\rho_0
\]
is supported in $\R^d\times\overline{B(0,M)}$, namely
\bq\label{eq:compact_support_xi}
\supp\mu^{(R)}_s \subset \R^d\times\overline{B(0,M)} .
\eq
In particular, for each $s\in[0,T]$, the restriction of $\mu^{(R)}_s$ to $K\times\R^d$ is a finite positive Radon measure supported in the compact set $K\times\overline{B(0,M)}$, and hence the family $\{\mu^{(R)}_s\}_{R>0}$ is bounded in $\calM(K\times\overline{B(0,M)})$.

Using \eqref{eq:compact_support_xi} and the fact that $\mu^{(R)}_s$ are probability measures, we may extract a sequence $R_n\to\infty$ and a limit
\[
\mu^\ast \in L^\infty\bigl(0,T;\calM(K\times\overline{B(0,M)})\bigr), 
\]
such that $\mu^{(R_n)}\rightharpoonup\mu^\ast$ weakly-$\ast$ in $L^\infty (0,T;\calM(K\times\overline{B(0,M)}) )$. That is, for every test function
$\psi\in L^1(0,T;C(K\times\overline{B(0,M)}))$, we have
\[%\bq\label{eq:mu_conv}
\int_0^T \intrr \psi(s,z,\xi)\,\mu^{(R_n)}_s(\dz,\rd\xi)\,\ds \to \int_0^T \intrr \psi(s,z,\xi)\,\mu^\ast_s(\dz,\rd\xi)\,\ds.
\]%\eq

For a.e.  $s\in(0,T)$, let $\rho^\ast_s$ denote the $z$--marginal of $\mu^\ast_s$. By the disintegration theorem (Theorem \ref{thm:disint}), there exists a $\rho^\ast_s$--a.e.  uniquely determined family of probability measures $\{\tilde\nu^\ast_{s,z}\}_{z\in K}\subset\calP(\R^d)$ such that
\[
\mu^\ast_s(\dz,\rd\xi)=\rho^\ast_s(\dz)\,\tilde\nu^\ast_{s,z}(\rd\xi).
\]
%We define the associated barycentric velocity by
%\[
%u^\ast_s(z):=\intr \xi\,\tilde\nu^\ast_{s,z}(\rd\xi), \quad \rho^\ast_s\text{-a.e.  } z\in K.
%\]

\medskip
\noindent\textbf{Step 3. Equicontinuity in time and compactness.} We show that the family $\{\mu^{(R)}\}_{R>0}$ is equicontinuous in time in the dual space $W^{-1,\infty}(K\times\overline{B(0,M)})$.

Let $\psi\in C^1(K\times\overline{B(0,M)})$. From the kinetic formulation \eqref{eq:kinetic_weak_R}, we obtain
\begin{align*}
\frac{\rd}{\rd s}\intrr\psi(z,\xi)\,\mu^{(R)}_s(\dz,\rd\xi) &= \intrr \xi\cdot\nabla_z\psi(z,\xi)\,\mu^{(R)}_s(\dz,\rd\xi)\cr
&\quad + \intrr F_p[\mu^{(R)}_s](z,\xi)\cdot\nabla_\xi\psi(z,\xi)\,\mu^{(R)}_s(\dz,\rd\xi).
\end{align*}
Since $|\xi|\le M$ on the support of $\mu^{(R)}_s$ and the alignment force 
$F_p[\mu^{(R)}_s]$ is uniformly bounded on $K\times\overline{B(0,M)}$, we deduce
\[
\lt|\frac{\rd}{\rd s}\intrr\psi(z,\xi)\,\mu^{(R)}_s(\dz,\rd\xi) \rt|
\le C\|\psi\|_{W^{1,\infty}(K\times\overline{B(0,M)})},
\]
with a constant $C >0$ independent of $R$.  Hence $\{\mu^{(R)}\}_{R>0}$ is equicontinuous in  $W^{-1,\infty}(K\times\overline{B(0,M)})$.

Combining this equicontinuity with the weak-$\ast$ compactness obtained in Step 2,
an Arzel\'a--Ascoli argument yields, up to extraction of a subsequence,
\[
\mu^{(R_n)}_s \rightharpoonup \mu^\ast_s
\quad\text{in }\calM(K\times\overline{B(0,M)})
\quad\text{for every } s\in[0,T].
\]

\medskip
\noindent\textbf{Step 4. Collapse of the velocity marginal.}
For each $R>0$ and $t\in[0,T]$, let
\[
\lambda^{(R)}_t := (\pi_\xi)_\#\mu^{(R)}_t = (v_{R+t})_\#\rho_0
\in \calP(\overline{B(0,M)})
\]
be the velocity marginal of $\mu^{(R)}_t$, where $\pi_\xi(z,\xi):=\xi$ denotes the projection onto the velocity variable.

Since $\pi_\xi$ is continuous and $\mu^{(R_n)}_t \rightharpoonup \mu^\ast_t$ narrowly in
$\calM(K\times\overline{B(0,M)})$ for every $t\in[0,T]$, it follows that
\[
\lambda^{(R_n)}_t \rightharpoonup \lambda^\ast_t := (\pi_\xi)_\#\mu^\ast_t
\quad\text{narrowly in }\calP(\overline{B(0,M)})
\]
for every $t\in[0,T]$.

We now show that $\lambda^\ast_t$ is in fact a Dirac mass concentrated at the conserved mean velocity. Since $\lambda^{(R_n)}_t=(v_{R_n+t})_\#\rho_0$, every two points in $\supp\lambda^{(R_n)}_t$
are of the form $v_{R_n+t}(x)$ and $v_{R_n+t}(y)$ for some $x,y\in\supp\rho_0$.
Hence
\[
\diam(\supp\lambda^{(R_n)}_t) \le \rd_v(R_n+t).
\]
Therefore, for every fixed $t\in[0,T]$,
\[
\diam(\supp\lambda^{(R_n)}_t)\to0 \quad\text{as }n\to\infty.
\]

Next we use conservation of total momentum to identify the center of concentration. By Lemma \ref{lem:apr}, 
the total momentum is conserved along the Lagrangian dynamics:
\[
\intr v_s(x)\,\rho_0(\dx)=\intr u_0(x)\,\rho_0(\dx)=:\bar u \quad\text{for all }s\ge0.
\]
Equivalently,
\[
\intr \xi\,\lambda^{(R_n)}_t(\rd\xi)=\bar u \quad\text{for every }n\in\N,\ t\in[0,T].
\]

We claim that
\[
\lambda^{(R_n)}_t \rightharpoonup \delta_{\bar u} \quad\text{narrowly in }\calP(\overline{B(0,M)})
\]
for every fixed $t\in[0,T]$.
To prove this, let $g\in C(\overline{B(0,M)})$ be arbitrary.
Since $\bar u$ is the barycenter of $\lambda^{(R_n)}_t$, and the support of $\lambda^{(R_n)}_t$
has diameter tending to zero, the whole support must concentrate around $\bar u$.
Indeed, for any $\xi\in\supp\lambda^{(R_n)}_t$,
\[
|\xi-\bar u|
=
\lt|\xi-\intr \xi'\,\lambda^{(R_n)}_t(\rd\xi')\rt|
\le
\intr |\xi-\xi'|\,\lambda^{(R_n)}_t(\rd\xi')
\le
\diam(\supp\lambda^{(R_n)}_t).
\]
Taking the supremum over $\xi\in\supp\lambda^{(R_n)}_t$, we obtain
\[
\supp\lambda^{(R_n)}_t \subset
\overline{B\bigl(\bar u,\diam(\supp\lambda^{(R_n)}_t)\bigr)}.
\]
Hence,
\[
\lt| \intr g(\xi)\,\lambda^{(R_n)}_t(\rd\xi)-g(\bar u) \rt|
\le
\sup_{|\xi-\bar u|\le \diam(\supp\lambda^{(R_n)}_t)}
|g(\xi)-g(\bar u)|.
\]
Since $g$ is uniformly continuous on the compact set $\overline{B(0,M)}$ and
$\diam(\supp\lambda^{(R_n)}_t)\to0$, the right-hand side converges to $0$.
Therefore,
\[
\intr g(\xi)\,\lambda^{(R_n)}_t(\rd\xi)\to g(\bar u),
\]
which proves the claim.

By the uniqueness of the narrow limit, we conclude that
\[
\lambda^\ast_t=\delta_{\bar u} \quad\text{for every }t\in[0,T].
\]

Finally, since $\lambda^\ast_t=(\pi_\xi)_\#\mu^\ast_t$ is the $\xi$-marginal of $\mu^\ast_t$
and is equal to the Dirac mass $\delta_{\bar u}$, the measure $\mu^\ast_t$ must be concentrated on the slice
$\R^d\times\{\bar u\}$. Hence, we have
\[
\mu^\ast_t(\rd z,\rd\xi)=\rho^\ast_t(\rd z)\,\delta_{\bar u}(\rd\xi)
\quad\text{for every }t\in[0,T].
\]

\medskip
\noindent\textbf{Step 5. Passage to the limit and identification of the limit dynamics.}
Let $\varphi\in C_c^\infty((0,T)\times\R^d)$ be arbitrary, and let
$K\subset\R^d$ be a compact set such that
\[
\supp\varphi\subset (0,T)\times K.
\]
Since $\mu^{(R_n)}$ is a weak solution of the kinetic equation, testing against
functions depending only on $(s,z)$ yields
\[
\int_0^T\intrr
\lt(-\partial_s\varphi(s,z)-\xi\cdot\nabla_z\varphi(s,z)\rt) \mu^{(R_n)}_s(\rd z,\rd\xi)\,\rd s=0.
\]
Indeed, the force term disappears because $\varphi$ is independent of $\xi$.

Now the function
\[
(s,z,\xi)\mapsto -\partial_s\varphi(s,z)-\xi\cdot\nabla_z\varphi(s,z)
\]
belongs to $L^1(0,T;C(K\times\overline{B(0,M)}))$.
Hence, by the weak-$\ast$ convergence obtained in Step 2, we may pass to the limit:
\[
\int_0^T\intrr \lt(-\partial_s\varphi(s,z)-\xi\cdot\nabla_z\varphi(s,z)\rt) \mu^\ast_s(\rd z,\rd\xi)\,\rd s=0.
\]
Using the monokinetic form from Step 4,
\[
\mu^\ast_s(\rd z,\rd\xi)=\rho^\ast_s(\rd z)\,\delta_{\bar u}(\rd\xi),
\]
we obtain
\[
\int_0^T\intr \lt(-\partial_s\varphi(s,z)-\bar u\cdot\nabla_z\varphi(s,z)\rt) \rho^\ast_s(\rd z)\,\rd s=0.
\]
Therefore, we have
\[
\partial_t\rho^\ast + \nabla_z \cdot (\rho^\ast \bar u)=0 \quad\text{in }\calD'((0,T)\times\R^d).
\]

Settting
\[
\rho^\ast_0 := (\pi_z)_\#\mu^\ast_0 \in \calP(\R^d),
\]
the above continuity equation with constant velocity $\bar u$ yields
\[
\rho^\ast_t=(z\mapsto z+t\bar u)_\#\rho^\ast_0 \quad\text{for all }t\in[0,T].
\]
As $T>0$ was arbitrary, the same representation holds for all $t\ge0$.
This completes the proof.
\end{proof}

\end{document}
